% docs/vocab/vocab-macros.tex -- macros used by the shared vocabulary entries.
%
% Every definition is guarded with \providecommand so this file can be \input
% from the BODY of a host paper whose preamble already defines the same names:
%   HAaNE II (docs/paper-ii/paper.tex) defines \R, \vx, \vone, \KM identically
%   (\KM = \mathcal{M}); HAaNE I (docs/Final-twist/paper/preamble.tex) defines
%   \R identically but \KM = \textsc{KM} (a text abbreviation, unused in its
%   body) and writes the matrix as \M = \mathrm{M}; appendix_vocab_I.tex
%   therefore re-binds \KM, \vx and \vone inside a group. Do not add
%   \newcommand / \renewcommand / \def here.
% Requirements on the host: amsmath + amssymb (both papers load them).
%
% The entry macro: title, then the four fixed fields. Inside the entry
% paragraph a line may break after any underscore of a typewriter path (test
% and function names run to 100 characters, wider than a line, and cmtt does
% not hyphenate) and TeX may stretch inter-word space (\sloppy) instead of
% running into the margin; both settings are scoped to the entry.
\providecommand{\vocabusbreak}{\vocaborigus\allowbreak}
\providecommand{\vocabentry}[5]{\par\medskip\noindent\begingroup\sloppy\let\vocaborigus\_\let\_\vocabusbreak\textbf{#1.}\ \emph{Definition.} #2\ \emph{Algebraic reading.} #3\ \emph{Computed here.} #4\ \emph{Validated here.} #5\par\endgroup}
% Symbols the entries use.
\providecommand{\R}{\mathbb{R}}
\providecommand{\vx}{\boldsymbol{x}}
\providecommand{\vone}{\boldsymbol{1}}
\providecommand{\KM}{\mathcal{M}}
% Cross-reference to another entry, by its title.
\providecommand{\vocabsee}[1]{(see \emph{#1})}
